\documentclass[10pt,a4paper]{amsart}
\usepackage{amsmath,amssymb,amsthm,graphicx}
\usepackage{verbatim}
\newenvironment{reference}{\comment}{\endcomment}
%\newenvironment{reference}{}{}
\newenvironment{slogan}{\comment}{\endcomment}
\newenvironment{history}{\comment}{\endcomment}

% For commutative diagrams we use Xy-pic
\usepackage[all]{xy}

% We use 2cell for 2-commutative diagrams.
\xyoption{2cell}
% Theorem environments.
%
\theoremstyle{plain}
\newtheorem{theorem}[subsection]{Theorem}
\newtheorem{proposition}[subsection]{Proposition}
\newtheorem{lemma}[subsection]{Lemma}

\theoremstyle{definition}
\newtheorem{definition}[subsection]{Definition}
\newtheorem{example}[subsection]{Example}
\newtheorem{exercise}[subsection]{Exercise}
\newtheorem{situation}[subsection]{Situation}

\theoremstyle{remark}
\newtheorem{remark}[subsection]{Remark}
\newtheorem{remarks}[subsection]{Remarks}

\numberwithin{equation}{subsection}

% Macros
%
\def\lim{\mathop{\mathrm{lim}}\nolimits}
\def\colim{\mathop{\mathrm{colim}}\nolimits}
\def\Spec{\mathop{\mathrm{Spec}}}
\def\Hom{\mathop{\mathrm{Hom}}\nolimits}
\def\Ext{\mathop{\mathrm{Ext}}\nolimits}
\def\SheafHom{\mathop{\mathcal{H}\!\mathit{om}}\nolimits}
\def\SheafExt{\mathop{\mathcal{E}\!\mathit{xt}}\nolimits}
\def\Sch{\mathit{Sch}}
\def\Mor{\mathop{\mathrm{Mor}}\nolimits}
\def\Ob{\mathop{\mathrm{Ob}}\nolimits}
\def\Sh{\mathop{\mathit{Sh}}\nolimits}
\def\NL{\mathop{N\!L}\nolimits}
\def\CH{\mathop{\mathrm{CH}}\nolimits}
\def\proetale{{pro\text{-}\acute{e}tale}}
\def\etale{{\acute{e}tale}}
\def\QCoh{\mathit{QCoh}}
\def\Ker{\mathop{\mathrm{Ker}}}
\def\Im{\mathop{\mathrm{Im}}}
\def\Coker{\mathop{\mathrm{Coker}}}
\def\Coim{\mathop{\mathrm{Coim}}}

% Boxtimes
%
\DeclareMathSymbol{\boxtimes}{\mathbin}{AMSa}{"02}

%
% Macros for moduli stacks/spaces
%
\def\QCohstack{\mathcal{QC}\!\mathit{oh}}
\def\Cohstack{\mathcal{C}\!\mathit{oh}}
\def\Spacesstack{\mathcal{S}\!\mathit{paces}}
\def\Quotfunctor{\mathrm{Quot}}
\def\Hilbfunctor{\mathrm{Hilb}}
\def\Curvesstack{\mathcal{C}\!\mathit{urves}}
\def\Polarizedstack{\mathcal{P}\!\mathit{olarized}}
\def\Complexesstack{\mathcal{C}\!\mathit{omplexes}}
% \Pic is the operator that assigns to X its picard group, usage \Pic(X)
% \Picardstack_{X/B} denotes the Picard stack of X over B
% \Picardfunctor_{X/B} denotes the Picard functor of X over B
\def\Pic{\mathop{\mathrm{Pic}}\nolimits}
\def\Picardstack{\mathcal{P}\!\mathit{ic}}
\def\Picardfunctor{\mathrm{Pic}}
\def\Deformationcategory{\mathcal{D}\!\mathit{ef}}
\usepackage{cleveref}
\crefname{originalreference}{original reference}{original references}
\Crefname{originalreference}{Original reference}{Original references}
\makeatletter
\newlabel{spaces-morphisms-lemma-fppf-open}{{042S}{source}{}{sourcepointer-1}{}}
\newlabel{spaces-morphisms-lemma-fppf-open@cref}{{[originalreference][0][]042S}{[1][0][]source}{}{}{}}
\makeatother
\crefname{originalreference}{original reference}{original references}
\Crefname{originalreference}{Original reference}{Original references}
\makeatletter
\newlabel{spaces-properties-lemma-open-subspaces}{{03BZ}{source}{}{sourcepointer-2}{}}
\newlabel{spaces-properties-lemma-open-subspaces@cref}{{[originalreference][0][]03BZ}{[1][0][]source}{}{}{}}
\makeatother
\crefname{originallemma}{original lemma}{original lemmas}
\Crefname{originallemma}{Original lemma}{Original lemmas}
\makeatletter
\newlabel{lemma-pullback-property-local-target}{{06EM}{source}{}{sourcepointer-3}{}}
\newlabel{lemma-pullback-property-local-target@cref}{{[originallemma][0][]06EM}{[1][0][]source}{}{}{}}
\makeatother
\title{Largest open of the base where the restriction has a \$\textbackslash{}tau\$-local property}
\author{The Stacks Project Authors}
\date{stacks-project:spaces-descent.tex:local-snapshot-20261012; GNU FDL 1.2 or later}
\begin{document}
\maketitle
\noindent\textit{Editorial note.} All delivered mathematical text is the author's own, copied verbatim from the pinned Stacks Project chapter file spaces-descent.tex; no statement or proof step was rewritten, reordered or removed.\par\medskip
\noindent\textit{Editorial note.} The counted span cites one Morphisms of Spaces lemma, one Properties of Spaces lemma and the same-chapter pullback lemma, recorded in assembly.external\_references with canonical tags as navigation pointers.\par\medskip
\noindent\textit{Allowed standard prerequisites (editorial).}\begin{itemize}
\item Openness of fppf-type morphisms on points. Stacks Project, Morphisms of Spaces, Lemma spaces-morphisms-lemma-fppf-open. Shows \$W\_\{set\}\$ is open.
\item Open subspaces of algebraic spaces. Stacks Project, Properties of Spaces, Lemma spaces-properties-lemma-open-subspaces. Upgrades the open subset of points to an open subspace.
\item Base-change stability of \$\textbackslash{}tau\$-local properties. Stacks Project, Descent on Spaces, Lemma lemma-pullback-property-local-target. Supplies \$W(f') \textbackslash{}supset g\textasciicircum{}\{-1\}(W(f))\$.
\end{itemize}\medskip
\ifdefined\hypertarget\hypertarget{sourcepointer-1}{}\fi
\noindent\textit{Source reference.} Original reference 042S: sources/stacks-chow/tags.txt, source lines 5339--5339. Cross-chapter standard prerequisite cited by the author inside the counted span; it is used as a pointer and its own proof is not part of this item.\par\medskip
\ifdefined\hypertarget\hypertarget{sourcepointer-2}{}\fi
\noindent\textit{Source reference.} Original reference 03BZ: sources/stacks-chow/tags.txt, source lines 4366--4366. Cross-chapter standard prerequisite cited by the author inside the counted span; it is used as a pointer and its own proof is not part of this item.\par\medskip
\ifdefined\hypertarget\hypertarget{sourcepointer-3}{}\fi
\noindent\textit{Source reference.} Original lemma 06EM: sources/stacks-chow/tags.txt, source lines 8221--8221. Cross-chapter standard prerequisite cited by the author inside the counted span; it is used as a pointer and its own proof is not part of this item.\par\medskip
\begin{lemma}
\label{lemma-largest-open-of-the-base}
Let $S$ be a scheme.
Let $\tau \in \{fppf, syntomic, smooth, \etale\}$.
Let $\mathcal{P}$ be a property of morphisms of algebraic spaces over $S$
which is $\tau$ local on the target. For any morphism of algebraic spaces
$f : X \to Y$ over $S$ there exists a largest open subspace
$W(f) \subset Y$ such that the restriction $X_{W(f)} \to W(f)$ has
$\mathcal{P}$. Moreover,
\begin{enumerate}
\item if $g : Y' \to Y$ is a morphism of algebraic spaces which is
flat and locally of finite presentation, syntomic, smooth, or \'etale
and the base change $f' : X_{Y'} \to Y'$ has $\mathcal{P}$, then
$g$ factors through $W(f)$,
\item if $g : Y' \to Y$ is flat and locally of finite presentation,
syntomic, smooth, or \'etale, then $W(f') = g^{-1}(W(f))$, and
\item if $\{g_i : Y_i \to Y\}$ is a $\tau$-covering, then
$g_i^{-1}(W(f)) = W(f_i)$, where $f_i$ is the base change of $f$
by $Y_i \to Y$.
\end{enumerate}
\end{lemma}

\begin{proof}
Consider the union $W_{set} \subset |Y|$ of the images
$g(|Y'|) \subset |Y|$ of morphisms $g : Y' \to Y$ with the properties:
\begin{enumerate}
\item $g$ is flat and locally of finite presentation, syntomic,
smooth, or \'etale, and
\item the base change $Y' \times_{g, Y} X \to Y'$ has property
$\mathcal{P}$.
\end{enumerate}
Since such a morphism $g$ is open (see
Morphisms of Spaces, Lemma \ref{spaces-morphisms-lemma-fppf-open})
we see that $W_{set}$ is an open subset of $|Y|$. Denote $W \subset Y$
the open subspace whose underlying set of points is $W_{set}$, see
Properties of Spaces, Lemma \ref{spaces-properties-lemma-open-subspaces}.
Since $\mathcal{P}$ is local in the $\tau$ topology the restriction
$X_W \to W$ has property $\mathcal{P}$ because we are given a covering
$\{Y' \to W\}$ of $W$ such that the pullbacks have $\mathcal{P}$.
This proves the existence and proves that $W(f)$ has property (1).
To see property (2) note that $W(f') \supset g^{-1}(W(f))$ because
$\mathcal{P}$ is stable under base change by flat and locally of finite
presentation, syntomic, smooth, or \'etale morphisms, see
Lemma \ref{lemma-pullback-property-local-target}.
On the other hand, if $Y'' \subset Y'$ is an open such that
$X_{Y''} \to Y''$ has property $\mathcal{P}$, then $Y'' \to Y$ factors
through $W$ by construction, i.e., $Y'' \subset g^{-1}(W(f))$. This
proves (2). Assertion (3) follows from (2) because each morphism
$Y_i \to Y$ is flat and locally of finite presentation, syntomic,
smooth, or \'etale by our definition of a $\tau$-covering.
\end{proof}

\end{document}
